\scanpage{164}
\mspage{80}
% archive sheet number: 163

\diagram{\begin{tikzcd}[ampersand replacement=\&]
\unsure{\zeta|\xi} \arrow[r,shift left=.55ex] \arrow[r,shift right=.55ex] \& \varphi_\alpha(i),
\end{tikzcd}}

existe une flèche $i\overset{u}{\unsure{-|\longrightarrow}}j$ dans $I$ telle que $\varphi_\alpha(u)$ égalise la double flèche.

\uline{Dém} Si $\varphi_\alpha$ est cofinal, on en \unsure{conclut|déduit} par définition que $\varinjlim\deleted{\unsure{\psi|\varphi}}\varphi_\alpha\longrightarrow\varinjlim\psi\varphi_\alpha$ est iso, or cette flèche n'est autre que $\alpha$. Inversement, supposons que cette flèche soit un iso, donc $C_{/F}\xleftarrow{\sim}C_{/\varinjlim\deleted{\unsure{I}} x_i}\unsure{)}$, et il reste à voir que le foncteur can.

\[
I\longrightarrow C_{/\varinjlim_{(C^\wedge)}x_i}
\]

est cofinal, i.e. satisfait F1) et F2). La vérification est $\pm$ \unsure{tautologique}, en utilisant le fait que pour tt \deleted{\unsure{$\zeta$}}$y$ \deleted{\unsure{$\in$}} \inserted{$\leftarrow$} dans $C$ \deleted{(et \unsure{notamment un quelconque} \unsure{des termes} $x_j$, $j\in\operatorname{Ob} I$),}

\[
\text{on a}\qquad \operatorname{Hom}(y,\varinjlim_{C^\wedge}x_i)\xleftarrow{\sim}\varinjlim_I{}_{\mathrm{Ens}}\operatorname{Hom}(y,x_i),
\]

et du ``calcul'' des limites inductives filtrantes dans Ens. \unsure{$\mathsection$|$\square$}

\uline{Corollaire} 4.4.3. \inserted{Soit $\pi$ un cardinal \unsure{infini}.} Pour que $F$ soit dans $\operatorname{Ind}^{\unsure{+|'}}_{\pi}(C)$\newline
\begin{tabular}{@{}|p{.91\linewidth}@{}}
(image essentielle de $\operatorname{Ind}_{\pi}(C)$ dans $C^\wedge$) il f. et s. qu'il existe une catégorie ordonnée gr. de devant $\pi$ $I$, et un foncteur cofinal $I\longrightarrow C_{/F}$.
\end{tabular}

Il \deleted{faut donc} nous faut donc donner une condition, sur une catégorie $\unsure{J|\mathcal J}$ $(=C_{/F})$, pour qu'il existe une catégorie ordonnée $I$ \deleted{\unsure{grande}} devant $\pi$, et un foncteur cofinal $I\longrightarrow J$.

